\section{The fixed-margin Gaussian endpoint}
\label{sec:zero-endpoint}
We now prove the confidence lower bound for the exact $\gamma=0$
posterior. The full event $A_S$ from Section~\ref{sec:zero-conditioning}
uses the parity labels $y_i=f(X_i)$. Write $\mathbb E_S^*$ for expectation
under $G\mid A_S$, and let $\mathcal D=\mathcal D(S)$.
Its relative entropy from $N(0,K)$ is exactly $\mathcal D$.

\begin{lemma}[Gaussian quadratic entropy]\label{lem:zero-quadratic}
Suppose $Z$ is a linear image of a centered Gaussian vector under a
reference law, with covariance $T\succeq0$, trace $t$, and largest
eigenvalue $\ell>0$. Under any law of relative entropy at most $D_0$
from that reference law,
\[
 \mathbb E\|Z\|^2\le2t+4\ell D_0.
\]
\end{lemma}
\begin{proof}
The entropy variational inequality states that
$\mathbb E_Q V\le\operatorname{KL}(Q\Vert P)+\log\mathbb E_Pe^V$.
It follows immediately from nonnegativity of the relative entropy
from $e^VdP/\mathbb E_Pe^V$, first for bounded $V$ and then by
truncation. At $V=\|Z\|^2/(4\ell)$, the Gaussian moment generating
function gives
\[
 \log\mathbb E_Pe^V
 =-\frac12\sum_j\log\left(1-\frac{\lambda_j(T)}{2\ell}\right)
 \le\frac{t}{2\ell},
\]
using $-\log(1-x)\le2x$ for $0\le x\le1/2$. The claim follows.
The zero-covariance case is immediate.
\end{proof}

For any odd cutoff $r$, apply this lemma to the high-degree projection
$G_{>r}$ as a vector scaled by $M^{-1/2}$, and to its restrictions at
the $n$ training locations scaled by $n^{-1/2}$. The resulting exact
bounds are
\begin{align}
 \mathbb E_S^*\|G_{>r}\|_2^2
 &\le2\tau_r+4\lambda_{r+1}\mathcal D,
 \label{eq:zero-tail-population}\\
 \mathbb E_S^*\mathbb P_nG_{>r}^2
 &\le2\tau_r+\frac{4\|Q_{r,SS}\|\mathcal D}{n}.
 \label{eq:zero-tail-empirical}
\end{align}
Here $\|g\|_2^2=\mathbb E_xg(x)^2$ and
$\mathbb P_ng=n^{-1}\sum_i g(X_i)$. The normalized empirical covariance
has trace exactly $\tau_r$, even if locations repeat. Neither formula
asserts independence under the posterior, or log-concavity of the
globally conditioned law.

\begin{theorem}[Endpoint confidence lower bound]\label{thm:zero-endpoint}
For fixed $L\ge2$ and $a>0$, there are $c>0$ and $d_0$ such that, for
every $d\ge d_0$ and each integer $0\le n\le e^{cd}$, with probability
at least $3/4$ over independent uniform parity training inputs,
\[
 R_S:=\mathbb E_S^*\mathbb P_x\{f(x)G(x)<0\}\ge\frac38.
\]
Consequently the fraction of inputs whose correct-label posterior
probability is at most $2/3$ is at least $1/16$.
\end{theorem}
\begin{proof}
We first consider $1\le n\le\exp[(\log(d+1))^3]$.
On the exponentially high-probability event of
Lemma~\ref{lem:zero-denominator},
$\mathcal D\le C(d^2+nd^{1+\beta})$.
Let $s$ be odd with $s\sim d^{1/3}$ and put $p=G_{\le s}$.
The spectral estimates give $\tau_s=o(1)$ and
\[
 \mathcal D\lambda_{s+1}
 \le \frac{\operatorname{poly}(d)\exp[(\log(d+1))^3]}
                {\binom{d+1}{s+2}}=o(1),
\]
uniformly in this range of sample sizes. Thus
\eqref{eq:zero-tail-population} implies
$\mathbb E_S^*\|G-p\|_2^2=o(1)$.
Since $|G(x)|>a$ at every input, the signs of $G$ and $p$ can disagree
only where $|G-p|\ge a$, using any fixed sign convention at $p=0$.
Their posterior-averaged disagreement fraction is therefore $o(1)$.
Lemma~\ref{lem:ptf-parity} gives parity error $1/2-o(1)$ for every
degree-$s$ polynomial threshold, so $R_S\ge1/2-o(1)>3/8$.

For the remaining range, fix $\delta=1/20$. Choose a fixed sufficiently
large sampling constant $A$ and use the maximal odd cutoff of
Lemma~\ref{lem:adaptive-cutoff}:
\[
 A D_r\log(2D_r)\le n<A D_{r+2}\log(2D_{r+2}),
 \qquad \tfrac14(\log d)^2\le r\le\alpha d.
\]
Here $\alpha<1/256$ is a small positive constant to be chosen below;
the exponential constant $c$ is then chosen small enough to enforce
this upper bound. Write
\[
 \rho=r/d,\qquad H_\rho=\log(e/\rho),\qquad
 \tau=\tau_r,\quad\ell=\lambda_{r+1},\qquad
 B_r=C_\delta\bigl(1+\sqrt{A H_\rho/\rho}+A H_\rho\bigr).
\]
Using $r/d$ instead of $r/(d+1)$ in the adaptive lemma changes only
fixed constants.
With probability at least $1-2\delta-o(1)$, the denominator,
one-sided squared-loss sampling, and tail Gram events all hold:
\begin{align}
 \mathcal D/n&\le C d\rho^{-\beta},
 \label{eq:zero-normalized-cost}\\
 \mathbb P_n(fP)_-^2&\ge\tfrac14\|P\|_2^2
       \quad\hbox{for every degree-$r$ polynomial $P$},
 \label{eq:zero-coercivity}\\
 \|Q_{r,SS}\|&\le B_r\tau.
 \label{eq:zero-gram}
\end{align}
For the first inequality use \eqref{eq:zero-adaptive-cost},
$n\ge A D_r\log(2D_r)$, and $\tau\asymp r^{-\beta}$.
Indeed $d^2/n=o(1)$ here and
$D_r\log(r+1)/n\le C$, while the main term is
$d^{1+\beta}\tau\le C d\rho^{-\beta}$.
The other two inequalities follow from
\eqref{eq:wrong-square-coercivity} and
Lemma~\ref{lem:adaptive-cutoff}, respectively. Fix a sample on these
events for the rest of the argument.

Take $u=r+2\lceil10\log(d+1)\rceil$ and decompose the actual posterior
output into orthogonal Walsh bands,
\[
 G=p+b+t,\qquad p=G_{\le r},\quad b=G_{(r,u]},\quad t=G_{>u}.
\]
Every posterior draw has $f(X_i)G(X_i)>0$. Therefore, pointwise at
the training inputs, $(fp)_-\le|G_{>r}|$.
Equations~\eqref{eq:zero-coercivity}--\eqref{eq:zero-gram} and
\eqref{eq:zero-tail-empirical} give
\begin{equation}\label{eq:zero-head}
 \mathbb E_S^*\|p\|_2^2
 \le8\tau+16 B_r\tau\mathcal D/n.
\end{equation}
The middle band has prior trace at most $\tau$ and largest normalized
eigenvalue at most $\ell$. Lemma~\ref{lem:zero-quadratic} gives
\[
 \mathbb E_S^*\|b\|_2^2\le2\tau+4\ell\mathcal D.
\]
Orthogonality, together with the shell estimates
$n\ell/\tau\le CAH_\rho$ and
$n\lambda_{u+1}/\tau\le d^{-10}$ from
Lemma~\ref{lem:adaptive-cutoff}, now implies
\begin{align}
 \mathbb E_S^*\|p+b\|_2^2
 &\le10\tau+C\frac{\mathcal D\tau}{n}(B_r+AH_\rho),
 \label{eq:zero-low-norm}\\
 \mathbb E_S^*\|t\|_2^2
 &\le2\tau_u+4d^{-10}\frac{\mathcal D\tau}{n}.
 \label{eq:zero-final-tail}
\end{align}
No additional empirical sampling event at $u$ is needed.

Set $A_d=b_0d^{(1-\beta)/2}$, with the fixed constant $b_0>0$ of
Lemma~\ref{lem:zero-amplitude}. That lemma applies because
$\mathcal D\le e^{.045d}$ on the denominator event.
Its conclusion and \eqref{eq:zero-final-tail} yield, uniformly over
the sample sizes in question,
\begin{align}
 \mathbb E_S^*\mathbb P_x\{|G(x)|\le A_d\}&\le C/d,
 \label{eq:zero-large-outputs}\\
 \mathbb E_S^*\|t\|_2^2/A_d^2
 &\le C\{(d\rho^\beta)^{-1}+d^{-10}\rho^{-2\beta}\}=o(1).
 \label{eq:zero-tail-relative}
\end{align}
The last assertion uses $\rho\ge(\log d)^2/(4d)$ and $\beta\le1/2$.

Let $\operatorname{clip}(z)=\max(-1,\min(1,z))$.
The exact clipped-parity identity in Lemma~\ref{lem:clipped-parity},
applied to the degree-$u$ polynomial $p+b$, gives
\[
 \left|\mathbb E_x f(x)\operatorname{clip}((p+b)(x)/A_d)\right|
 \le\frac{u}{dA_d}\|p+b\|_2.
\]
Clipping is $1$-Lipschitz, and
$|\sign G-\operatorname{clip}(G/A_d)|\le\mathbf1_{\{|G|\le A_d\}}$.
Consequently Cauchy--Schwarz and the preceding bounds show
\begin{align}
 \left|\mathbb E_S^*\mathbb E_x f(x)\sign G(x)\right|
 &\le\frac{u}{dA_d}\sqrt{\mathbb E_S^*\|p+b\|_2^2}
       +\frac{\sqrt{\mathbb E_S^*\|t\|_2^2}}{A_d}+C/d.
 \label{eq:zero-correlation}
\end{align}
The second and third terms tend to zero uniformly. Since $u/r=1+o(1)$,
\eqref{eq:zero-normalized-cost} and \eqref{eq:zero-low-norm} bound
the square of the first term by
\begin{equation}\label{eq:zero-rho-powers}
 C\rho^{2-2\beta}
       \bigl(1+\rho^{-1/2}\sqrt{H_\rho}+H_\rho\bigr)+o(1).
\end{equation}
All constants here are independent of $d,n,\rho$; they may depend on
the fixed $L,a,A,\delta$.
The worst power is $\rho^{3/2-2\beta}\sqrt{H_\rho}$, which tends to
zero as $\rho\downarrow0$ for every $\beta\le1/2$.
Choose $\alpha$ small enough and then $d$ large enough that the right
side of \eqref{eq:zero-correlation} is at most $1/4$.
Thus
\[
 R_S=\frac{1-\mathbb E_S^*\mathbb E_x f(x)\sign G(x)}2\ge\frac38.
\]
The good sample event has probability at least
$1-2\delta-o(1)>3/4$. At $n=0$, global sign symmetry gives correct-label
probability $1/2$ at every input and the same risk conclusion directly.

Finally put $q_S(x)=\mathbb P\{f(x)G(x)>0\mid A_S\}$ and
$F_S=\mathbb P_x\{q_S(x)\le2/3\}$. Then
\[
 R_S=\mathbb E_x(1-q_S(x))
 \le\frac{1-F_S}{3}+F_S,
 \qquad F_S\ge\frac{3R_S-1}{2}\ge\frac1{16}.
\]
This conclusion concerns the actual posterior probability, without
identifying it with the sign of a posterior mean.
\end{proof}
